%% APPENDIX D -- the false-alarm expectation, the control ensemble, and the
%% execution-block quality criterion.  Owner: r11-chain.  Carries
%% \label{app:falsealarm} (04_method.tex, 08_backmatter.tex) and
%% \label{sec:blockquality} (04_method.tex).
\section{The false-alarm expectation and the control ensemble}
\label{app:falsealarm}

The paper's central claim is that the number of unattributed crossings is the
number chance predicts. This appendix gives that expectation, what it is
conditioned on, and what the control ensemble calibrating it can support.

\subsection{What the expectation is conditioned on}
%% (subsection label retired: nothing cites it)

The expectation is built from the control statistics themselves, which are
the only empirical null the data provide, and it imposes the condition an
event must actually satisfy: to reach the trigger, and to fall outside the
attribution windows. Summed window by window over the \NWinA{} Class~A
windows, each window's own effective cell count against its own realised null
scale, \RsevChanceTrigA{} are expected to trigger by chance; the attribution
windows occupy \RsevChanceOccPct{} per cent of the searched band; and the
product is \RsevChanceExp{} expected unattributed Class~A crossings.

That figure assumes the star's own tail matches its controls', and
Appendix~\ref{app:radius} measures that it does not: a no-signal stellar
position reaches the top of its own ensemble \HoTailFactor{} times as often
as exchangeability would give, with a bootstrap interval
\HoTailLo--\HoTailHi{} clustered over the reserved blocks. Folding that in,
the expectation is \textbf{about \ChnChanceMid{} unattributed Class~A
crossings, in a range \ChnChanceLo--\ChnChanceHi, against \ChnUnattrA{}
observed}. The range comes almost entirely from the width on that factor and
is far wider than the difference between the two numbers, so what the
comparison establishes is that the observed population needs no explanation
beyond chance --- not that chance predicts it to within one crossing.

The classes divide as \ChnXrossA{} Class~A crossings and \ChnXrossB{}
Class~B, of which \ChnAttrA{} and \ChnAttrB{} respectively are attributed; so
every Class~B crossing is unattributed, and all of them come from the one
execution block of \S\ref{sec:blockquality}.

\emph{One conditioning is reported only as a diagnostic.} An expectation
conditioned on the rank predicts a different quantity --- the number of
windows in which a no-signal position would rank first, \ClustStageMean{}
against the \ClustStageNObs{} observed --- and that is a statement about the
pipeline and not about the crossings. It also explains why every crossing
outranking all its controls is Class~A: only \StageNullEligible{} of the
\StageNullAllWin{} windows retaining a control vector have an ensemble
maximum reaching $T_{\rm trig}=\StageNullTrig$ at all, and
\StageNullEligibleA{} of those are Class~A.

\emph{One trade is not available to us, and it should be stated plainly.} An
excess of this size needs only a small unmodelled tail in the control
ensemble, and Appendix~\ref{app:radius} shows that one exists --- but the
same noise model sets the injection recovery fractions, so a tail large
enough to remove the excess would make the quoted completeness optimistic by
about the same factor. Neither correction has been applied.

\subsection{Execution-block quality}
\label{sec:blockquality}

The comparison with the control positions assumes each block's ensemble
behaves like noise, and one block makes that assumption plainly false: every
one of its \NCtrl{} probes exceeds the trigger, the smallest at
\BqFailCtrlMin. A criterion for the condition was therefore fixed and applied
to every block, rather than recognised on the block that prompted it.

The statistic is the one this paper already uses as a diagnostic: the median
of a window's control ensemble, which is what separates a field with a source
in it from a field that is no null at all. That median is a maximum over the
window's own channel$\times$drift grid, so its level rises with the cell
count; it is therefore scaled by the median of the same quantity over the
windows of its own search class, and one threshold then means the same thing
in both. A block's statistic is the median of that ratio over its windows,
and \textbf{a block fails when it exceeds \BqFactor{} the scale of its
class} --- the factor already carried by the window-level data-quality flag,
fixed before that flag was applied to anything.

Applied to all \BqNBlocks{} blocks of the primary census, where the median
block sits at \BqQMed, the criterion fails \BqNFail: \BqFailBlock{} toward
\BqFailStar, at \BqFailQ{} where the next worst block is \BqNextQ. Nothing
else is close, and the unscaled form of the rule fails the same
\BqRawNFail{} block. \textbf{The complete Class~A results are identical with
the failing block and without it} --- \BqWinA{} windows, \BqXrossA{}
crossings, \BqUnattrA{} unattributed, \BqRankA{} outranking all \NCtrl{}
controls --- because its \BqFailNWin{} windows are all coarse-channel ones.
What it does carry is \BqFailNCross{} of the survey's \NCross{} crossings,
the entire Class~B crossing population.

\subsection{The control ensemble: its calibration, and its resolution}
%% (subsection label retired: nothing cites it)

The ring geometry is set in \S\ref{sec:statistic}; its calibration and its
resolution are the subject here. Under that geometry the stellar ranks are
consistent with the uniform distribution exchangeability implies, in both
array strata --- \VerTwmAfterN{} 12\,m windows and \VerAcaAfterN{} ACA
windows, the latter with median rank \VerAcaAfterMed{} against the one half
expected --- over the \VerNTested{} of \NWindows{} windows that retain a
control vector. The ACA windows are measured on the geometry of each
observation's own antenna table and baselines rather than on a 12\,m beam.
Two cautions keep this from being a proof of exchangeability: a shift of less
than about \VerKsPowerAca{} in the rank distribution could not have been seen
at all in samples of this size, so the strata are consistent with uniform
rather than demonstrably uniform; and the diagnostic most sensitive to
contamination of the inner controls, the excess of the innermost control bin
over the outer annulus, is \VerAcaAfterExc{} with an interval
\VerAcaAfterExcLo{} to \VerAcaAfterExcHi{} that spans zero.

\emph{The rank is not a significance test, and the reason is structural.}
Ranking the star against \NCtrl{} control positions drawn from the same field
is a good way to decide which windows deserve examination, but
$1/(\NCtrl+1)=\RankFloorVal$ is the \emph{resolution} of the ensemble --- the
smallest rank it can express --- and not the probability that noise produced
the window. Against a survey-wide Bonferroni scale
%
\begin{equation}
\alpha_{\rm survey}=\frac{0.05}{\NWindows}=\BonfScaleCalc,
\qquad
\frac{\RankFloorVal}{\alpha_{\rm survey}}=\BonfRankRatio,
\label{eq:bonf}
\end{equation}
%
that floor is \BonfRankRatio{} times too coarse, and correlation widens it:
neighbouring controls are correlated below the synthesised-beam scale, so a
window's annulus carries only $N_{\rm eff}\simeq30$--$43$ independent spatial
trials and the ensemble resolves no better than \RankEffLo--\RankEffHi.
Drawing more controls does not rescue it, the primary beam bounding how many
independent positions a single pointing contains. \textbf{The floor of
Eq.~\ref{eq:bonf} is the design's central limitation: no single window in
this survey can be significant on the rank alone, which is why confirmation
here requires independent recurrence.}

\subsection{The trials budget}

A window holds \CellsWinMin--\CellsWinMax{} channel$\times$drift cells
(median \CellsWinMed), the survey \CellsTotal{} in total, and they are not
independent trials: channel smoothing and a drift grid oversampled two to
one together over-count them by \FaOverCount. The control ensembles
calibrate that budget without a Gaussian assumption. The same arithmetic
predicts \NCtrlWinExp{} of the \NWindows{} windows carrying a control above
$T_{\rm trig}$ and \NCtrlWinObs{} do, agreeing to \CtrlWinAgreePct{} per
cent; corrected for the over-count, the observed control maxima sit
$+\FaExcessFine$ and $+\FaExcessCoarse$ per cent above prediction in the
fine and coarse classes independently. Departures of that size are why the
operative statement is the empirical rate and not a Gaussian budget.

\emph{The residual on-star excess is not zero, and it is reported as an open
systematic rather than explained away.} Correlation cannot move the expected
\emph{number} above the trigger, only its clumping, so the excess over an
independent-Gaussian budget divides in two: the cells within one crossing
window, a factor \CellsPerCrossRatio{} that the frequency-axis over-count
predicts because one real excursion is counted about three times, and the
number of crossing windows, a factor \CrossWinRatio{} with a one-sided
Poisson probability of \CrossWinPoisson. The residual runs one way and
carries the same sign at the controls, so a residual excess of chance
crossings at the tens-of-per-cent level cannot be excluded.

\emph{One stratum departs from uniformity and it is the one that matters.}
The fine-channel windows' ranks are not uniform over $n=\NFine$ windows:
\NStageOneFine{} of them contain an on-star peak above their own control
maximum against \ExpStageOneFine{} expected, while the coarse stratum, three
times larger and predicting \ExpStageOneCoarse{} such windows, contains
\NStageOneCoarse. A coarse channel dilutes a narrow circumstellar line, so
the class that can resolve one is the class that finds it --- and by the same
reasoning a genuine unresolved carrier would appear in Class~A first. The
departure survives removal of every fine window carrying a stellar crossing,
so fine-window rank probabilities are optimistic for the class as a whole and
not only where a crossing sits.

\emph{What is left is bounded, and the bounds are null.} Treating each
control as a pseudo-star and ranking it against the other \NCtrl-1 gives
\NPseudo{} ranks that are uniform to the resolution the ensemble can
express, and the real stellar ranks agree with them at $\pv{\StarPseudoP}$.
Any systematic radial suppression of the controls is below
\RingRhoBoundPct{} per cent in the mean, and the star's own residual
variance equals the controls' to better than a per cent in the rank-passing
windows. Scrambling each window so that nothing coherent survives along a
drift track carries \LNNWindows{} windows past the rank floor and changes no
disposition; it fails only on HD~48370 Band~6, whose extended emission and
imaging residuals are coherent across integrations and are destroyed by the
scramble itself, so no scrambled-noise probability is quoted there.

%% ------------------------------------------------------------------
%% NOTE (r11-chain, round 11).  THE QUESTION THE PREVIOUS NOTE ASKED IS
%%   ANSWERED AND THE ANSWER IS THE OTHER WAY ROUND.  This appendix now
%%   leads with the expectation for the chain actually applied -- trigger
%%   plus attribution, no rank -- carries the measured tail factor on it, and
%%   quotes the result as ONE range (\ChnChanceMid, \ChnChanceLo-\ChnChanceHi
%%   against \ChnUnattrA).  The rank-conditioned expectation (\ClustStageMean
%%   against \ClustStageNObs) is kept as a pipeline diagnostic and is no
%%   longer reachable as "the" expectation.  \RsevChanceFracLoc and
%%   \RsevChanceObsAUnattr are no longer cited here: the first belonged to a
%%   localisation clause the chain no longer has, and the second is the
%%   RELEASED CATALOGUE's unattributed count, which disagrees with the
%%   ledger's (see INTEGRATION_QUEUE entry 2).
%% REQUEST (appendix-triage -> integrator):
%%   Floats deleted from this appendix and moved to the data release:
%%   tab:ctrlsep (control separations), tab:excess (the cell-excess
%%   decomposition), tab:arraysplit (the array split) and fig:ctrldiag
%%   (control diagnostics).  None is cited from the body.  tab_ctrlsep_v400,
%%   tab_falsealarm_v399 and the control_diagnostics figure are now
%%   uninputted; retire_macros will strip their macros.
%% ------------------------------------------------------------------

%% ★ tab:ledger MOVED HERE from Sec. 5 in the v4.12 length pass: the
%% crossing-by-crossing ledger belongs beside the false-alarm accounting that
%% explains how many of its rows are expected.  0.72 pp off the main text.
\begin{table*}
\centering
\caption{\textbf{The crossing ledger, from which every disposition in this
paper can be reproduced.} Every one of the \NCross{} threshold crossings:
the star, its execution block, the ALMA band, the crossing frequency as the
correlator delivered it, the same frequency in the star's rest frame, the
search statistic $T_\star$, the largest statistic $T_{\rm ctrl}$ among the
window's \NCtrl{} control positions, the nearest masked transition and the
stellar-frame offset $\Delta v_\star$ from it, and then the recurrence
test: $n_{\rm rep}$ covering repeat blocks spanning $\Delta t$ days, the
statistic $T_{\rm pers}$ a carrier persisting at the discovery flux would
have produced at the predicted stellar-frame cell, the statistic
$T_{\rm rep}$ measured there at the matched drift, and the exclusion.
$T_{\rm ctrl}$ is a diagnostic and gates nothing: it says which crossing was
examined first. The drift-maximised reading of the same cell is in
Table~\ref{tab:cand} for the two leading events and in the data release for
all of them. One execution block fails the quality criterion of
\S\ref{sec:bdblock}, and the two quantities that would inherit its
discovery amplitude --- $T_{\rm pers}$ and the exclusion --- are withheld
on its four rows rather than printed from data the paper calls defective.
A crossing is attributed where $|\Delta v_\star|\le\MaskHalfKms$\,km\,s$^{-1}$;
the transition named is the nearest in that frame, so far from every
transition it is the magnitude and not the name that carries information.
$\nu_\star$ follows from $\nu_{\rm topo}$ by the barycentric term of that
block's epoch and the star's systemic velocity and nothing else, so the last
column is reproducible from the two before it. Block identifiers omit the
constant \texttt{A002\_} prefix.}
\label{tab:ledger}
\tiny
%% GENERATED by ledger_v410.py -- do not hand-edit.
%% One row per threshold crossing, one epoch convention, stellar-frame attribution.
\begin{tabular}{@{}l@{\hspace{0.9pt}}l@{\hspace{0.9pt}}c@{\hspace{0.9pt}}r@{\hspace{0.9pt}}r@{\hspace{0.9pt}}r@{\hspace{0.9pt}}r@{\hspace{0.9pt}}l@{\hspace{0.9pt}}r@{\hspace{0.9pt}}r@{\hspace{0.9pt}}r@{\hspace{0.9pt}}r@{\hspace{0.9pt}}r@{\hspace{0.9pt}}r@{\hspace{0.9pt}}l@{}}
\hline
Star & block & B & $\nu_{\rm topo}$ (GHz) & $\nu_{\star}$ (GHz) & $T_\star$ & $T_{\rm ctrl}$ & line & $\Delta v_{\star}$ & $n_{\rm rep}$ & $\Delta t$ (d) & $T_{\rm pers}$ & $T_{\rm rep}$ & $\sigma_{\rm excl}$ & disposition \\
\hline
BD+05 1668 & Xc7fa6f\_X28a8 & 6 & 240.1016 & 240.1054 & 73.398 & 19.73 & CS($5$--$4$) & -5911.9 & 29 & 0.9--113 & n/a & +1.81 & n/a & unattributed; flagged \\
BD+05 1668 & Xc7fa6f\_X28a8 & 6 & 224.7422 & 224.7458 & 48.303 & 19.51 & H$_2$CO($3_{1,2}$--$2_{1,1}$) & -1264.5 & 29 & 0.9--113 & n/a & +1.62 & n/a & unattributed; flagged \\
BD+05 1668 & Xc7fa6f\_X28a8 & 6 & 226.7422 & 226.7458 & 41.355 & 18.40 & CN($2_{3/2}$--$1_{1/2}$) & +88.0 & 29 & 0.9--113 & n/a & +2.37 & n/a & unattributed; flagged \\
BD+05 1668 & Xc7fa6f\_X28a8 & 6 & 242.1953 & 242.1992 & 39.586 & 18.55 & CS($5$--$4$) & -3349.2 & 29 & 0.9--113 & n/a & +1.42 & n/a & unattributed; flagged \\
$\beta$~Pic & Xf5d76d\_Xeb8 & 3 & 115.2603 & 115.2711 & 30.765 & 7.34 & CO($1$--$0$) & -0.2 & --- & --- & --- & --- & --- & attributed \\
$\beta$~Pic & Xf5d76d\_Xe19 & 3 & 115.2603 & 115.2711 & 27.684 & 6.48 & CO($1$--$0$) & -0.2 & --- & --- & --- & --- & --- & attributed \\
HD 48370 & Xc26103\_X155a & 6 & 230.5117 & 230.5243 & 26.834 & 18.36 & CO($2$--$1$) & -17.8 & --- & --- & n/a & --- & n/a & attributed; flagged \\
$\beta$~Pic & Xf4df6f\_Xc7 & 3 & 115.2620 & 115.2711 & 17.907 & 5.76 & CO($1$--$0$) & -0.2 & --- & --- & --- & --- & --- & attributed \\
$\beta$~Pic & Xf5d76d\_Xcc5 & 3 & 115.2604 & 115.2711 & 17.288 & 5.52 & CO($1$--$0$) & -0.2 & --- & --- & --- & --- & --- & attributed \\
$\beta$~Pic & Xf5d76d\_X32f1 & 3 & 115.2606 & 115.2713 & 14.654 & 7.27 & CO($1$--$0$) & +0.4 & --- & --- & --- & --- & --- & attributed \\
$\beta$~Pic & Xd9668b\_X3a90 & 6 & 230.5161 & 230.5378 & 11.681 & 9.12 & CO($2$--$1$) & -0.3 & --- & --- & --- & --- & --- & attributed \\
$\beta$~Pic & X6f1341\_X1484 & 6 & 230.5284 & 230.5379 & 9.832 & 8.20 & CO($2$--$1$) & -0.1 & --- & --- & --- & --- & --- & attributed \\
$\beta$~Pic & X7116f1\_X1785 & 6 & 230.5272 & 230.5384 & 9.677 & 14.13 & CO($2$--$1$) & +0.5 & --- & --- & --- & --- & --- & attributed \\
$\beta$~Pic & Xd9668b\_Xa9df & 6 & 230.5160 & 230.5377 & 8.948 & 8.41 & CO($2$--$1$) & -0.4 & --- & --- & --- & --- & --- & attributed \\
$\beta$~Pic & Xa7a216\_X23e4 & 6 & 230.5273 & 230.5377 & 7.988 & 6.03 & CO($2$--$1$) & -0.4 & --- & --- & --- & --- & --- & attributed \\
61 Vir & Xc079b5\_X82f & 7 & 346.3138 & 346.3232 & 5.958 & 5.65 & SO($9_{8}$--$8_{7}$) & -177.6 & 3 & 0.1--3.0 & 9.7 & +1.65 & 9.3 & unattributed; no recurrence \\
HD 285968 & Xd98580\_X3bf6 & 6 & 230.5021 & 230.5447 & 5.956 & 8.12 & CO($2$--$1$) & +8.7 & --- & --- & --- & --- & --- & attributed \\
HD 285968 & Xdb7ab7\_X5b4e & 6 & 230.5116 & 230.5447 & 5.868 & 7.22 & CO($2$--$1$) & +8.7 & --- & --- & --- & --- & --- & attributed \\
HD 285968 & Xdb4b9a\_X108f & 6 & 230.5100 & 230.5447 & 5.835 & 7.94 & CO($2$--$1$) & +8.7 & --- & --- & --- & --- & --- & attributed \\
CP$-$72 2713 & Xff0235\_X4a6d & 7 & 344.2697 & 344.2965 & 5.809 & 5.68 & SO($8_{8}$--$7_{7}$) & -12.3 & 5 & 0.1--2.1 & 6.3 & +0.96 & 7.1 & unattributed; no recurrence \\
HD~139084~B & Xcd8029\_Xb6b0 & 7 & 345.8241 & 345.8241 & 5.763 & --- & CO($3$--$2$) & +24.4 & --- & --- & --- & --- & --- & attributed \\
HD 14055 & Xfff3d8\_Xd46 & 7 & 331.7699 & 331.7664 & 5.575 & 5.99 & $^{13}$CO($3$--$2$) & +1068.7 & 30 & 0.9--650 & 170.4 & +2.35 & 170.4 & unattributed; no recurrence \\
HD 23484 & X10b22d7\_X2445 & 6 & 230.7904 & 230.8006 & 5.520 & 5.55 & CO($2$--$1$) & +341.5 & 4 & 0.1--2.1 & 10.4 & +1.53 & 9.2 & unattributed; no recurrence \\
HD 161868 & Xff45a6\_X10ca8 & 7 & 345.6338 & 345.6491 & 5.481 & 6.16 & SO($2_{3}$--$2_{1}$) & -48.0 & 13 & 0.0--6.1 & 92.2 & +1.56 & 92.0 & attributed \\
HIP 10679 & Xa95c04\_X14cd & 6 & 220.4185 & 220.4035 & 5.456 & 6.76 & $^{13}$CO($2$--$1$) & +6.5 & --- & --- & --- & --- & --- & attributed \\
HD~14055 & X11adad7\_X25c5e & 7 & 345.4564 & 345.4303 & 5.448 & --- & SO($2_{3}$--$2_{1}$) & -237.8 & 30 & 0.1--653 & 14.0 & +2.48 & 14.0 & unattributed; no recurrence \\
HD 207129 & Xc19d6f\_X5706 & 6 & 230.4226 & 230.4070 & 5.433 & 5.52 & CO($2$--$1$) & -170.3 & 4 & 0.8--425 & 69.7 & +0.86 & 70.7 & unattributed; no recurrence \\
$\beta$~Pic & Xd9668b\_X3a90 & 6 & 241.7512 & 241.7739 & 5.420 & 5.81 & CS($5$--$4$) & -3869.7 & 1 & 1.0 & 5.7 & +0.72 & 5.0 & unattributed; no recurrence \\
HR~1010 & Xc6fa08\_X1053 & 6 & 230.5133 & 230.5287 & 5.410 & --- & CO($2$--$1$) & -12.0 & --- & --- & --- & --- & --- & attributed \\
ALMA J153702653-33192492 & Xff0235\_X8392 & 6 & 230.5222 & 230.5392 & 5.405 & 13.41 & CO($2$--$1$) & +1.6 & --- & --- & --- & --- & --- & attributed \\
HD 110411 & Xe45e29\_Xb7d4 & 6 & 219.4538 & 219.4357 & 5.393 & 5.85 & C$^{18}$O($2$--$1$) & -170.2 & 1 & 10 & 6.7 & +0.62 & 6.1 & unattributed; no recurrence \\
HD 31392 & X129754b\_X6946 & 6 & 231.1998 & 231.2231 & 5.289 & 5.81 & H($30\alpha$) & -876.3 & 8 & 0.1--41 & 15.7 & +1.74 & 16.0 & unattributed; no recurrence \\
HD 14055 & Xff99e1\_X24be0 & 7 & 331.7078 & 331.7035 & 5.277 & 5.69 & $^{13}$CO($3$--$2$) & +1011.6 & 30 & 0.9--652 & 252.2 & +2.14 & 252.0 & unattributed; no recurrence \\
AU Mic & Xa42f75\_X319 & 6 & 230.6864 & 230.6705 & 5.274 & 6.07 & CO($2$--$1$) & +172.3 & 2 & 310--455 & 6.4 & -0.09 & 6.5 & unattributed; no recurrence \\
HD 69830 & X12209f7\_X16edb & 7 & 322.4325 & 322.4444 & 5.259 & 5.42 & C$^{18}$O($3$--$2$) & -6268.5 & 5 & 0.1--29 & 8.9 & +1.69 & 9.4 & unattributed; no recurrence \\
ALMA J153702653-33192492 & Xfe62c1\_X398 & 6 & 216.1844 & 216.2036 & 5.228 & 6.02 & SiO($5$--$4$) & -1244.6 & 7 & 18--126 & 11.0 & +1.75 & 10.0 & unattributed; no recurrence \\
AU Mic & X899169\_X1838 & 6 & 230.8252 & 230.8294 & 5.221 & 5.85 & CO($2$--$1$) & +379.0 & 2 & 145--310 & 10.0 & +0.69 & 9.8 & unattributed; no recurrence \\
HD 14055 & Xff99e1\_X2c3bb & 7 & 345.5256 & 345.5216 & 5.186 & 5.66 & SO($2_{3}$--$2_{1}$) & -158.6 & 30 & 0.1--651 & 172.8 & +2.08 & 170.8 & unattributed; no recurrence \\
HD 14055 & Xfffde1\_X3835 & 7 & 345.0099 & 345.0070 & 5.177 & 6.56 & SO($2_{3}$--$2_{1}$) & -604.8 & 30 & 0.1--649 & 187.2 & +1.38 & 189.0 & unattributed; no recurrence \\
HD 14055 & Xff99e1\_X2b2ee & 7 & 344.7007 & 344.6965 & 5.175 & 5.92 & SO($8_{8}$--$7_{7}$) & +336.0 & 30 & 0.1--651 & 210.7 & +2.47 & 210.7 & unattributed; no recurrence \\
LHS 1140 & Xd9d7c7\_X6fd2 & 6 & 230.6355 & 230.6268 & 5.136 & 5.68 & CO($2$--$1$) & +115.5 & 4 & 5.0--27 & 7.3 & +1.03 & 6.9 & unattributed; no recurrence \\
HD 10647 & Xff294f\_Xaf70 & 7 & 331.6570 & 331.6943 & 5.120 & 5.91 & $^{13}$CO($3$--$2$) & +1003.3 & 10 & 0.1--592 & 125.8 & +1.33 & 125.4 & unattributed; no recurrence \\
CD$-$57 1054 & Xcf525f\_X5aa4 & 7 & 346.5099 & 346.5265 & 5.108 & 5.71 & SO($9_{8}$--$8_{7}$) & -1.7 & 1 & 39 & 6.0 & -0.27 & 6.3 & attributed \\
ALMA J153702653-33192492 & Xf8f6a9\_X10359 & 6 & 217.7292 & 217.7269 & 5.091 & 5.41 & H$_2$CO($3_{0,3}$--$2_{0,2}$) & -680.5 & 7 & 3.8--242 & 11.6 & +1.43 & 11.5 & unattributed; no recurrence \\
GJ 849 & Xd9668b\_X841a & 6 & 230.8769 & 230.8581 & 5.090 & 5.83 & CO($2$--$1$) & +416.3 & 5 & 0.1--51 & 9.3 & -0.02 & 11.0 & unattributed; no recurrence \\
ALMA J153702653-33192492 & Xfe62c1\_X398 & 6 & 218.8729 & 218.8923 & 5.090 & 5.56 & H$_2$CO($3_{2,1}$--$2_{2,0}$) & +181.2 & 7 & 18--126 & 10.5 & +1.52 & 11.5 & unattributed; no recurrence \\
HD 14055 & X1001fdc\_X29ae & 7 & 331.1559 & 331.1546 & 5.082 & 6.26 & $^{13}$CO($3$--$2$) & +513.8 & --- & --- & --- & --- & --- & unattributed; no spectrum \\
$\beta$~Pic & Xd9668b\_Xa9df & 6 & 241.8367 & 241.8594 & 5.079 & 5.99 & CS($5$--$4$) & -3765.1 & 1 & 1.0 & 4.8 & -0.63 & 5.4 & unattributed; no recurrence \\
ALMA J153702653-33192492 & Xf8f6a9\_X10359 & 6 & 218.6711 & 218.6688 & 5.048 & 5.77 & H$_2$CO($3_{2,1}$--$2_{2,0}$) & -125.1 & 7 & 3.8--242 & 11.6 & +1.18 & 11.6 & unattributed; no recurrence \\
HD 69830 & X12277ed\_X178c6 & 7 & 321.2005 & 321.2149 & 5.046 & 6.13 & C$^{18}$O($3$--$2$) & -7387.7 & 5 & 6.0--36 & 10.2 & +1.26 & 10.6 & unattributed; no recurrence \\
HR 1010 & Xc66ce1\_Xaa9 & 6 & 230.8348 & 230.8491 & 5.044 & 5.72 & CO($2$--$1$) & +404.6 & 7 & 1.1--14 & 10.3 & +0.54 & 11.8 & unattributed; no recurrence \\
61 Vir & Xc067f7\_X822 & 7 & 345.5599 & 345.5678 & 5.038 & 5.79 & SO($2_{3}$--$2_{1}$) & -118.5 & 3 & 2.0--3.0 & 9.3 & +0.60 & 8.8 & unattributed; no recurrence \\
HD 170773 & Xff0235\_X8c42 & 7 & 331.0101 & 331.0234 & 5.029 & 5.75 & $^{13}$CO($3$--$2$) & +394.8 & 4 & 0.9--113 & 74.8 & +2.25 & 73.3 & unattributed; no recurrence \\
ALMA J153702653-33192492 & Xd9d7c7\_Xbb64 & 7 & 347.2179 & 347.1890 & 5.009 & 5.68 & SiO($8$--$7$) & -122.2 & 1 & 0.1 & 5.0 & -2.06 & 7.1 & unattributed; no recurrence \\
HD~14055 & X1190de8\_Xd526 & 7 & 344.4676 & 344.4515 & 5.006 & --- & SO($8_{8}$--$7_{7}$) & +122.7 & 30 & 20--611 & 11.6 & +2.24 & 12.2 & unattributed; no recurrence \\
$\eta$~Crv & X122b6ff\_X1041e & 7 & 357.2815 & 357.2475 & 5.006 & 5.50 & HCO$^{+}$($4$--$3$) & +431.3 & 1 & 87 & 11.0 & -0.62 & 11.6 & unattributed; no recurrence \\
\hline
\end{tabular}

\end{table*}
